\section{Реализация алгоритма}

Полноценная реализация алгоритма со всеми оптимизациями требует значительного количества кода, поэтому в этом разделе описаны только основные подходы. Для простоты примеры кода приводятся на языке программирования \OCaml.

Предположим, что термы заданы следующим алгебраическим типом данных:
\begin{minted}{OCaml}
type term =
| Var of int
| Con of int * term array
\end{minted}


\subsection{Система уравнений}

По аналогии с треугольными подстановками~\cite{baader2001unification}, наши системы уравнений естественным образом представляются в <<треугольном>> виде. Для этого мы используем тип частичных отображений из номеров переменных в термы \mintinline{OCaml}|term M.t|:
\begin{minted}{OCaml}
module M = Map.Make(Int)
\end{minted}

Во всех операциях над системами уравнений, мы используем номера представляющих переменных вместо $\E$. Таким образом, \walk реализуется следующим образом:
\begin{minted}{OCaml}
let rec walk s x = match M.find x s with
| exception Not_found -> x, Var x
| Var x -> walk s x
| t -> x, t
\end{minted}

Простейшая реализация операции \union, соответственно:
\begin{minted}{OCaml}
let union s x y =
  let x, xt = walk s x in
  let y, yt = walk s y in
  if x == y then s, None
  else begin match xt, yt with
  | Var _, Var _ -> M.add x yt s, None
  | Var _, _ -> M.add x (Var y) s, None
  | _, Var _ -> M.add y (Var x) s, None
  | _, _ -> M.add x (Var y) s, Some (xt, yt)
  end
\end{minted}

На практике, при реализации \union стоит применять эвристику <<возраста переменных>>~\cite{mannila1986complexity}: по возможности связывать <<более свежие>> переменных (обладающие б\'{о}льшими номерами) с <<менее свежими>>. Обоснование данной эвристики описано в процитированной статье, но на практике она позволяет уменьшить глубину рекурсии \walk подобно тому, как это делает ранговая эвристика в структуре данных СНМ (система непересекающихся множеств~\cite{galler1964improved}). Примеры реализации с использованием этой оптимизации приведены в репозитории с анализом производительности алгоритма.


\subsection{Общая часть термов}

Простейшая реализация \commonpart повторяет определение:
\begin{minted}{OCaml}
let rec common_part x y = match x, y with
| Var x, _ -> Var x
| _, Var y -> Var y
| Con (f, ts1), Con (g, ts2)
when f == g && Array.length ts1 == Array.length ts2 ->
  Con (f, Array.map2 common_part ts1 ts2)
| _ -> raise No_common_part
\end{minted}

На практике, полезно, чтобы реализация \commonpart возвращала непосредственно один из своих аргументов всегда, когда это возможно. Это позволяет снизить затраты памяти и увеличить шанс срабатывания быстрой проверки по равенству ссылок (<<physical equality>>) при выполнении унификации.

Для этого можно рассмотреть следующую решётку из четырёх элементов (<<сторона>>, \textsf{Side}): $\top$, $\Leftarrow$, $\Rightarrow$ и $\bot$.
\begin{center}
    \begin{tikzpicture}
        \node (both) at (0, 1) {$\top$};
        \node (left) at (-1, 0) {$\Leftarrow$};
        \node (right) at (1, 0) {$\Rightarrow$};
        \node (aside) at (0, -1) {$\bot$};
        \draw[dotted] (left) -- (both) node[midway,sloped] {$<$};
        \draw[dotted] (right) -- (both) node[midway,sloped] {$>$};
        \draw[dotted] (aside) -- (left) node[midway,sloped] {$>$};
        \draw[dotted] (aside) -- (right) node[midway,sloped] {$<$};
        \node[anchor=south] at (both.north) {<<обе>>};
        \node[anchor=east] at (left.west) {<<слева>>};
        \node[anchor=west] at (right.east) {<<справа>>};
        \node[anchor=north] at (aside.south) {<<ни одна>>};
    \end{tikzpicture}
\end{center}

Тогда можно реализовать обобщённую функцию $\commonpart' : T \times T \rightarrow (T \times \textsf{Side})^?$, которая вторым возвращаемым параметром в случае успеха дополнительно сообщает, какие из аргументов подходят на роль результата. Нерекурсивные вызовы в реализации очевидны, а при рекурсивных вызовах требуется накапливать инфинум <<сторон>>:
\begin{itemize}
    \item Если накопленный инфинум -- $\bot$, функция вынуждена вернуть сконструированный терм, потому что никакой из аргументов не подходит на роль результата.
    \item Иначе -- отбросить сконструированный терм и вернуть тот аргумент, на который указывает накопленный инфинум.
\end{itemize}

Конкретные примеры реализации расположены в репозитории с анализом производительности.


\subsection{Алгоритм}

Наконец, реализация \unify почти полностью повторяет псевдокод:
\begin{minted}{OCaml}
let rec unify_vt s x yt =
  let x, xt = walk s x in
  begin match xt with
  | Var _ -> M.add x yt s
  | _ ->
    let t = common_part xt yt in
    unify (M.add x t s) xt yt
  end

and unify s xt yt =
  if xt == yt then s
  else match xt, yt with
  | Var x, Var y ->
    begin match union s x y with
    | s, Some (xt, yt) -> unify s xt yt
    | s, None -> s
    end
  | Var x, _ -> unify_vt s x yt
  | _, Var y -> unify_vt s y xt
  | Con (f, ts1), Con (g, ts2)
  when f == g && Array.length ts1 == Array.length ts2 ->
    let n = Array.length ts1 in
    let rec hlp i s =
      if i == n then s
      else hlp (i + 1) @@ unify s ts1.(i) ts2.(i)
    in
    hlp 0 s
  | _ -> raise No_common_part
\end{minted}
